\documentclass[a4paper,11pt]{article}
\usepackage[utf8]{inputenc}
\usepackage[T1]{fontenc}
\usepackage[french]{babel}
\usepackage[left=2cm,right=2cm,top=2cm,bottom=2cm]{geometry}
\usepackage{xcolor}
\usepackage{hyperref}
\usepackage{fontawesome5}
\usepackage{parskip}

% --- Couleurs Deloitte ---
\definecolor{primary}{RGB}{0, 105, 55}
\hypersetup{colorlinks=true, urlcolor=primary}

\begin{document}

% --- EXPÉDITEUR ---
\noindent
\begin{minipage}[t]{0.5\textwidth}
    \textbf{Loïc Aron MBASSI EWOLO}\\
    Élève-Ingénieur Bac+5 -- Double diplôme ENSEA\\[3pt]
    \faMapMarker\ Cergy, France\\
    \faPhone\ (+33) 7 59 52 33 98\\
    \faEnvelope\ \href{mailto:loicaron.mbassiewolo@ensea.fr}{loicaron.mbassiewolo@ensea.fr}
\end{minipage}
\hfill
% --- DATE ---
\begin{minipage}[t]{0.4\textwidth}
    \raggedleft
    Cergy, le 2 mars 2026
\end{minipage}

\vspace{1.2cm}

% --- DESTINATAIRE ---
\hfill
\begin{minipage}[t]{0.5\textwidth}
    \raggedleft
    \textbf{Deloitte Société d'Avocats}\\
    Équipe Gi3 IT\\
    Paris La Défense, France
\end{minipage}

\vspace{1cm}

\noindent\textbf{Objet :} Candidature -- Stage Consultant Informatique \& Financement de l'Innovation

\vspace{0.8cm}

Madame, Monsieur,

Analyser les travaux R\&D d'un client, identifier les verrous technologiques et rédiger la documentation technique associée : ce sont des compétences que je mobilise depuis trois ans, d'abord en recherche appliquée en IA (MILA, Google DeepMind), puis sur mes propres projets d'innovation logicielle. C'est ce profil scientifique et rédactionnel que je souhaite mettre au service de l'équipe Gi3 IT de Deloitte pour accompagner vos clients dans leurs démarches de Crédit d'Impôt Recherche.

Mon parcours de recherche m'a formé à l'analyse d'état de l'art et à l'identification de l'avancée scientifique par rapport aux connaissances existantes. Au MILA (Institut Québécois d'IA), j'ai étudié le phénomène de Grokking sur des réseaux de neurones : recherche bibliographique approfondie, reproduction d'expériences SOTA, analyse mathématique de la convergence et rédaction de rapports de synthèse. Le projet UROP sur l'analyse d'ECG, mené en collaboration avec des mentors Google DeepMind, m'a confronté aux mêmes exigences : situer les travaux par rapport à la littérature, documenter les challenges techniques rencontrés et les résultats obtenus. Ces expériences m'ont donné la rigueur analytique nécessaire pour évaluer si un projet répond aux critères d'éligibilité du CIR.

Au-delà de la recherche, ma couverture technique est large : IA générative (LLM, RAG, LangChain), développement logiciel (Python, PHP, Java, cloud/Docker), systèmes embarqués (C/C++, FPGA) et data science (TensorFlow, PyTorch). Cette polyvalence me permettra d'échanger avec des interlocuteurs variés, qu'il s'agisse de directeurs techniques dans l'édition logicielle, l'aéronautique ou l'énergie. Par ailleurs, mes +3 ans d'expérience en développement Full Stack avec des clients internationaux (Belgique, Canada, Québec) ont renforcé mes compétences relationnelles et ma capacité à comprendre rapidement des contextes métier différents. La rédaction d'articles techniques sur Medium et de documentation open source complète ce profil orienté communication écrite.

Je suis disponible dès avril 2026 pour un stage de 4 à 6 mois sur le site de Paris La Défense. L'opportunité de combiner expertise technique et conseil en financement de l'innovation, au contact de secteurs industriels variés, correspond à ma volonté de valoriser mon profil d'ingénieur au-delà du développement pur.

Je reste à votre disposition pour un entretien afin de discuter de ma candidature.

Je vous prie d'agréer, Madame, Monsieur, l'expression de mes salutations distinguées.

\vspace{0.8cm}

\textbf{Loïc Aron MBASSI EWOLO}\\
\textit{Élève-Ingénieur ENSEA / Polytechnique Yaoundé}

\end{document}
